\chapter*{Annexe A -- Extrait de code}
\addcontentsline{toc}{chapter}{Annexe A -- Extrait de code}
\label{annexe:code}

Cette annexe reproduit l'implémentation de la fonction d'encodage
d'une série de trafic horaire en niveaux A-F (section~\ref{sec:methodo},
chapitre~\ref{chap:etatart}), citée dans le corps du rapport
(chapitre~\ref{chap:conception}).

\begin{lstlisting}[language=Python, basicstyle=\small\ttfamily, frame=single, caption={Encodage vectorisé d'une série horaire en niveaux A-F}, label={lst:encode}]
LEVEL_THRESHOLDS = [0.43, 0.63, 1.00, 1.35, 2.10]
ALPHABET = "ABCDEF"

def encode_levels(values, thresholds=LEVEL_THRESHOLDS):
    """Encode le trafic horaire d'une grille en niveaux 0..5,
    relatifs a sa mediane."""
    median = values.median()
    if median == 0:
        return pd.Series(0, index=values.index)
    ratio = values / median
    return np.searchsorted(thresholds, ratio, side="right")
\end{lstlisting}

\clearpage
\chapter*{Annexe B -- Extrait de code (fenêtre glissante)}
\addcontentsline{toc}{chapter}{Annexe B -- Extrait de code (fenêtre glissante)}
\label{annexe:code-fenetre}

Cette annexe reproduit les deux fonctions au cœur de la détection par
fenêtre glissante (section~\ref{sec:fenetre-glissante},
chapitre~\ref{chap:conception})~: le calcul de la distance
d'édition sur une fenêtre de longueur fixe glissée heure par heure, et
le regroupement des alertes consécutives par une même anomalie
prolongée (ne conserver que le \og{}front montant\fg{}).

\begin{lstlisting}[language=Python, basicstyle=\small\ttfamily, frame=single, caption={Distance sur fenêtre glissante et regroupement des alertes}, label={lst:sliding}]
def sliding_distances(obs, ref, win, block_dates):
    """Distance pour chaque fenetre glissante ; timestamp = (date,
    heure) de la FIN de la fenetre."""
    base = block_dates[0]
    out_d, out_ts = [], []
    for i in range(0, len(obs) - win + 1):
        out_d.append(DamerauLevenshtein.distance(
            obs[i:i+win], ref[i:i+win]))
        end_pos = i + win - 1
        out_ts.append((base + pd.Timedelta(days=end_pos // 24),
                        end_pos % 24))
    return out_d, out_ts

def debounce(is_anomaly):
    """Ne garde que les fronts montants : une anomalie qui dure 24h
    compte pour 1 incident, pas pour 24 alertes distinctes."""
    onsets = []
    prev = False
    for i, v in enumerate(is_anomaly):
        if v and not prev:
            onsets.append(i)
        prev = v
    return onsets
\end{lstlisting}

\clearpage
\chapter*{Annexe C -- Liste des abréviations}
\addcontentsline{toc}{chapter}{Annexe C -- Liste des abréviations}
\label{annexe:abreviations}

\begin{table}[htbp]
  \centering
  \begin{tabular}{ll}
    \toprule
    \textbf{Abréviation} & \textbf{Signification} \\
    \midrule
    4G/5G  & Réseaux mobiles de quatrième et cinquième génération \\
    RAN    & \emph{Radio Access Network} (réseau d'accès radio) \\
    UE     & \emph{User Equipment} (terminal utilisateur) \\
    CDR    & \emph{Call Detail Record} (enregistrement détaillé d'activité) \\
    ARIMA  & \emph{AutoRegressive Integrated Moving Average} \\
    SVM    & \emph{Support Vector Machine} \\
    SAX    & \emph{Symbolic Aggregate approXimation} \\
    PAA    & \emph{Piecewise Aggregate Approximation} \\
    UCL / LCL & Limite de contrôle supérieure / inférieure (\emph{Upper/Lower Control Limit}) \\
    ODbL   & \emph{Open Database License} \\
    ADN    & Signature journalière de 24 lettres (par analogie avec DiNATrAX) \\
    \bottomrule
  \end{tabular}
\end{table}
